\chapter{Particle in a Box in Two and Three Dimensions}
\index{particle in a box!two-dimensional|textbf}
\index{particle in a box!three-dimensional|textbf}
\index{Kronecker product|textbf}
\index{tensor-product basis|textbf}

\label{ch:training-particle-box-higher-dimensions}

\chapterstatus{pedagogical extension of the one-dimensional box. The chapter
derives separable continuum solutions and sparse finite representations. Its
examples are synthetic mathematics, not material calculations, scientific
validation, or uncertainty quantification.}

\section{Learning goals}

After working through this chapter, a reader should be able to

\begin{enumerate}
  \item construct two- and three-dimensional box eigenfunctions from
        one-dimensional eigenfunctions;
  \item explain what a Kronecker product does to matrices, vectors, dimensions,
        and tensor-product coordinates;
  \item distinguish a Kronecker product from a Kronecker sum;
  \item assemble two- and three-dimensional finite-difference Hamiltonians with
        a declared flattening order;
  \item derive product eigenvectors and additive energies; and
  \item explain symmetry-induced degeneracy in square and cubic boxes.
\end{enumerate}

Chapter~\ref{ch:training-quantum-operators} solved the one-dimensional box.
The higher-dimensional rectangular box is the cleanest setting in which to see
how separate coordinate directions combine. The continuum construction uses
tensor products of state spaces. Its finite matrix representation uses
Kronecker products.

\section{From an interval to a rectangular box}
\index{rectangular box}
\index{Dirichlet boundary condition}

A two-dimensional rectangular box occupies
\begin{equation}
  \Omega_2=(0,L_x)\times(0,L_y).
\end{equation}
A three-dimensional rectangular box occupies
\begin{equation}
  \Omega_3=(0,L_x)\times(0,L_y)\times(0,L_z).
\end{equation}
The wavefunction vanishes whenever any coordinate reaches a boundary. In two
dimensions,
\begin{equation}
  \psi(0,y)=\psi(L_x,y)=0,
  \qquad
  \psi(x,0)=\psi(x,L_y)=0.
\end{equation}
The corresponding three-dimensional statement also imposes zero values on the
two faces normal to the \(z\) direction.

Inside either box the potential is zero. The Hamiltonian is the Dirichlet
realization of the kinetic operator
\begin{equation}
  \hat H_d=-\frac{\hbar^2}{2m}\nabla_d^2,
\end{equation}
where \(d=2\) or \(3\). Cartesian derivatives separate:
\begin{align}
  \nabla_2^2
  &=\frac{\partial^2}{\partial x^2}
    +\frac{\partial^2}{\partial y^2},\\
  \nabla_3^2
  &=\frac{\partial^2}{\partial x^2}
    +\frac{\partial^2}{\partial y^2}
    +\frac{\partial^2}{\partial z^2}.
\end{align}
Each term changes one coordinate while leaving the others untouched. That
simple sentence is exactly what the later identity factors in the Kronecker
representation will encode.

\section{What a Kronecker product actually does}
\index{Kronecker product!definition}

Let
\begin{equation}
  A\in\mathbb R^{m\times n},
  \qquad
  B\in\mathbb R^{p\times q}.
\end{equation}
The Kronecker product \(A\otimes B\) is the block matrix
\begin{equation}
  A\otimes B
  =
  \begin{bmatrix}
    a_{11}B & a_{12}B & \cdots & a_{1n}B\\
    a_{21}B & a_{22}B & \cdots & a_{2n}B\\
    \vdots & \vdots & \ddots & \vdots\\
    a_{m1}B & a_{m2}B & \cdots & a_{mn}B
  \end{bmatrix}
  \in\mathbb R^{mp\times nq}.
\end{equation}
Every scalar entry \(a_{ij}\) of \(A\) is replaced by a complete scaled copy
\(a_{ij}B\). Thus the Kronecker product is not elementwise multiplication, and
it is not ordinary matrix multiplication. It constructs one larger linear map
from two maps acting on separate coordinate factors.

\subsection{A literal block example}

Take
\begin{equation}
  A=
  \begin{bmatrix}
    1 & 2\\
    3 & 4
  \end{bmatrix},
  \qquad
  B=
  \begin{bmatrix}
    0 & 5\\
    6 & 7
  \end{bmatrix}.
\end{equation}
Then
\begin{align}
  A\otimes B
  &=
  \begin{bmatrix}
    1B & 2B\\
    3B & 4B
  \end{bmatrix}\\
  &=
  \begin{bmatrix}
    0 & 5 & 0 & 10\\
    6 & 7 & 12 & 14\\
    0 & 15 & 0 & 20\\
    18 & 21 & 24 & 28
  \end{bmatrix}.
\end{align}
Both inputs are \(2\times2\), but their Kronecker product is \(4\times4\).
The block structure retains which factor supplied each action.

\subsection{Kronecker products of vectors}

For column vectors
\begin{equation}
  \mathbf a=
  \begin{bmatrix}a_1\\a_2\end{bmatrix},
  \qquad
  \mathbf b=
  \begin{bmatrix}b_1\\b_2\\b_3\end{bmatrix},
\end{equation}
the Kronecker product is
\begin{equation}
  \mathbf a\otimes\mathbf b
  =
  \begin{bmatrix}
    a_1\mathbf b\\
    a_2\mathbf b
  \end{bmatrix}
  =
  \begin{bmatrix}
    a_1b_1\\a_1b_2\\a_1b_3\\
    a_2b_1\\a_2b_2\\a_2b_3
  \end{bmatrix}
  \in\mathbb R^6.
\end{equation}
The product lists every pair of components in a specific order. Reversing the
factors generally changes that order:
\begin{equation}
  \mathbf a\otimes\mathbf b
  \neq
  \mathbf b\otimes\mathbf a
\end{equation}
as coordinate columns, even though a permutation matrix relates them.

The central action identity is
\begin{equation}
  (A\otimes B)(\mathbf x\otimes\mathbf y)
  =(A\mathbf x)\otimes(B\mathbf y),
\end{equation}
whenever the dimensions are compatible. The larger matrix applies \(A\) to the
first factor and \(B\) to the second factor. This identity is why Kronecker
products are natural for separable multidimensional operators.

More generally, compatible products obey the mixed-product rule
\begin{equation}
  (A\otimes B)(C\otimes D)
  =(AC)\otimes(BD).
\end{equation}
They also obey
\begin{equation}
  (A\otimes B)^{\mathsf T}=A^{\mathsf T}\otimes B^{\mathsf T},
\end{equation}
and, for complex matrices,
\begin{equation}
  (A\otimes B)^{\dagger}=A^{\dagger}\otimes B^{\dagger}.
\end{equation}
Consequently, the Kronecker product of Hermitian matrices is Hermitian.

\subsection{Identity factors mean ``leave this coordinate alone''}

Let \(A_x\) act on an \(N_x\)-component \(x\)-coordinate vector, and let
\(I_y\) be the \(N_y\times N_y\) identity. Then
\begin{equation}
  I_y\otimes A_x
\end{equation}
contains \(N_y\) copies of \(A_x\) on its block diagonal. It applies the same
\(x\)-direction action independently at every \(y\) index. The identity factor
means that the \(y\) coordinate is carried through unchanged.

Conversely,
\begin{equation}
  A_y\otimes I_x
\end{equation}
uses \(A_y\) to couple different \(y\) slices while \(I_x\) preserves the
\(x\) index inside each slice. These statements depend on the declared
flattening order, which we now make explicit.

Suppose samples are stored in an array \(\Psi_{j_y,j_x}\) and flattened in C
order, so that \(x\) varies fastest:
\begin{equation}
  \operatorname{vec}_C(\Psi)
  =
  \begin{bmatrix}
    \Psi_{1,1} & \cdots & \Psi_{1,N_x} &
    \Psi_{2,1} & \cdots & \Psi_{N_y,N_x}
  \end{bmatrix}^{\mathsf T}.
\end{equation}
With this ordering, \(I_y\otimes A_x\) acts along \(x\), and
\(A_y\otimes I_x\) acts along \(y\). A different flattening convention changes
the displayed matrices and requires a corresponding permutation. The physical
operator does not choose an array ordering for us.

\subsection{Kronecker product versus Kronecker sum}
\index{Kronecker sum}

For square matrices \(A\in\mathbb R^{n\times n}\) and
\(B\in\mathbb R^{p\times p}\), one common convention defines the Kronecker sum
by
\begin{equation}
  A\oplus B
  =I_p\otimes A+B\otimes I_n.
\end{equation}
The symbol \(\oplus\) here does not mean a direct sum of block-diagonal spaces.
It names the particular operator that applies \(A\) in one factor plus \(B\)
in the other. Because conventions can reverse factor order, a computational
record should state the expanded formula rather than rely on the symbol alone.

A separable Hamiltonian is represented by a Kronecker \emph{sum}, not by the
single product \(H_y\otimes H_x\). The sum encodes
\begin{equation}
  \text{energy from the }x\text{ direction}
  +
  \text{energy from the }y\text{ direction}.
\end{equation}

\section{The two-dimensional box}
\index{separation of variables}

The two-dimensional Hamiltonian is
\begin{equation}
  \hat H_2
  =-\frac{\hbar^2}{2m}
   \left(
     \frac{\partial^2}{\partial x^2}
     +\frac{\partial^2}{\partial y^2}
   \right).
\end{equation}
Seek a separated eigenfunction
\begin{equation}
  \psi(x,y)=X(x)Y(y).
\end{equation}
Substitution into \(\hat H_2\psi=E\psi\) gives
\begin{equation}
  -\frac{\hbar^2}{2m}
  \left(X''(x)Y(y)+X(x)Y''(y)\right)
  =EX(x)Y(y).
\end{equation}
Away from zeros of the factors, divide by \(X(x)Y(y)\):
\begin{equation}
  -\frac{\hbar^2}{2m}\frac{X''(x)}{X(x)}
  -\frac{\hbar^2}{2m}\frac{Y''(y)}{Y(y)}
  =E.
\end{equation}
The first term depends only on \(x\), the second only on \(y\), and their sum
is constant. Assign one-dimensional energies \(E_{n_x}^{(x)}\) and
\(E_{n_y}^{(y)}\). Each factor obeys the one-dimensional box equation and its
own boundary conditions.

The normalized product eigenfunctions are
\begin{equation}
  \psi_{n_x,n_y}(x,y)
  =\frac{2}{\sqrt{L_xL_y}}
   \sin\!\left(\frac{n_x\pi x}{L_x}\right)
   \sin\!\left(\frac{n_y\pi y}{L_y}\right),
\end{equation}
where
\begin{equation}
  n_x,n_y=1,2,3,\ldots.
\end{equation}
Their energies add:
\begin{equation}
  E_{n_x,n_y}
  =\frac{\hbar^2\pi^2}{2m}
   \left(
     \frac{n_x^2}{L_x^2}
     +\frac{n_y^2}{L_y^2}
   \right).
\end{equation}
The product states form an orthonormal basis of
\(L^2(0,L_x)\otimes L^2(0,L_y)\), naturally identified with
\(L^2(\Omega_2)\).

\subsection{Degeneracy in a square box}

If \(L_x=L_y=L\), interchanging \(n_x\) and \(n_y\) leaves the energy
unchanged:
\begin{equation}
  E_{n_x,n_y}
  =\frac{\hbar^2\pi^2}{2mL^2}
   \left(n_x^2+n_y^2\right)
  =E_{n_y,n_x}.
\end{equation}
For \(n_x\neq n_y\), the two product states are linearly independent, so this
exchange produces at least a two-dimensional degenerate eigenspace. For
example,
\begin{equation}
  E_{1,2}=E_{2,1}.
\end{equation}
The Hamiltonian fixes the eigenspace spanned by
\(\psi_{1,2}\) and \(\psi_{2,1}\), but any orthonormal rotation inside that
space is also a valid eigenbasis. A rectangular box with \(L_x\neq L_y\)
generally breaks this exchange degeneracy, though accidental coincidences may
still occur for particular length ratios and quantum numbers.

\section{Finite representation of the two-dimensional box}
\index{finite difference!two-dimensional box}
\index{Kronecker sum!two-dimensional Hamiltonian}

Let
\begin{equation}
  H_x\in\mathbb R^{N_x\times N_x},
  \qquad
  H_y\in\mathbb R^{N_y\times N_y}
\end{equation}
be the one-dimensional Dirichlet finite-difference Hamiltonians from
Chapter~\ref{ch:training-quantum-operators}. Flatten the grid with \(x\)
varying fastest. The two-dimensional represented Hamiltonian is
\begin{equation}
  H_{2,h}
  =I_{N_y}\otimes H_x
   +H_y\otimes I_{N_x}
  \in\mathbb R^{N_xN_y\times N_xN_y}.
\end{equation}
The first term applies \(H_x\) independently within every \(y\) slice. The
second term couples the slices according to \(H_y\) while preserving each
\(x\) index.

Suppose
\begin{equation}
  H_x\mathbf v_{n_x}^{(x)}
  =\varepsilon_{n_x}^{(x)}\mathbf v_{n_x}^{(x)},
  \qquad
  H_y\mathbf v_{n_y}^{(y)}
  =\varepsilon_{n_y}^{(y)}\mathbf v_{n_y}^{(y)}.
\end{equation}
Then the product vector
\begin{equation}
  \mathbf v_{n_y}^{(y)}\otimes\mathbf v_{n_x}^{(x)}
  \in\mathbb R^{N_xN_y}
\end{equation}
is an eigenvector of \(H_{2,h}\). Apply each term separately:
\begin{align}
  (I_{N_y}\otimes H_x)
  (\mathbf v_{n_y}^{(y)}\otimes\mathbf v_{n_x}^{(x)})
  &=\mathbf v_{n_y}^{(y)}\otimes
    (\varepsilon_{n_x}^{(x)}\mathbf v_{n_x}^{(x)}),\\
  (H_y\otimes I_{N_x})
  (\mathbf v_{n_y}^{(y)}\otimes\mathbf v_{n_x}^{(x)})
  &=(\varepsilon_{n_y}^{(y)}\mathbf v_{n_y}^{(y)})
    \otimes\mathbf v_{n_x}^{(x)}.
\end{align}
Adding the results gives
\begin{equation}
  H_{2,h}
  (\mathbf v_{n_y}^{(y)}\otimes\mathbf v_{n_x}^{(x)})
  =
  (\varepsilon_{n_x}^{(x)}+\varepsilon_{n_y}^{(y)})
  (\mathbf v_{n_y}^{(y)}\otimes\mathbf v_{n_x}^{(x)}).
\end{equation}
Thus product eigenvectors accompany additive energies in the discrete problem
just as they do in the continuum problem.

\subsection{A \(2\times2\) interior grid}

For equal grid spacings, suppress the common kinetic prefactor and use the
one-dimensional core
\begin{equation}
  K=
  \begin{bmatrix}
    2 & -1\\
    -1 & 2
  \end{bmatrix}.
\end{equation}
With two interior points along each axis,
\begin{align}
  K_2
  &=I_2\otimes K+K\otimes I_2\\
  &=
  \begin{bmatrix}
    4 & -1 & -1 & 0\\
    -1 & 4 & 0 & -1\\
    -1 & 0 & 4 & -1\\
    0 & -1 & -1 & 4
  \end{bmatrix}.
\end{align}
The four rows correspond, in one declared C-order convention, to
\begin{equation}
  (y_1,x_1),\ (y_1,x_2),\ (y_2,x_1),\ (y_2,x_2).
\end{equation}
Entries from \(I_2\otimes K\) connect points with the same \(y\) index and
neighboring \(x\) indices. Entries from \(K\otimes I_2\) connect points with
the same \(x\) index and neighboring \(y\) indices. The zero corner entries
show that this axial Laplacian does not directly connect diagonal neighbors.

\section{The three-dimensional box}

The normalized continuum eigenfunctions are
\begin{equation}
  \psi_{n_x,n_y,n_z}(x,y,z)
  =\sqrt{\frac{8}{L_xL_yL_z}}
   \sin\!\left(\frac{n_x\pi x}{L_x}\right)
   \sin\!\left(\frac{n_y\pi y}{L_y}\right)
   \sin\!\left(\frac{n_z\pi z}{L_z}\right),
\end{equation}
with energies
\begin{equation}
  E_{n_x,n_y,n_z}
  =\frac{\hbar^2\pi^2}{2m}
   \left(
     \frac{n_x^2}{L_x^2}
     +\frac{n_y^2}{L_y^2}
     +\frac{n_z^2}{L_z^2}
   \right).
\end{equation}
For a cube, \(L_x=L_y=L_z\), permutations of
\((n_x,n_y,n_z)\) have the same energy. The states associated with
\((1,1,2)\), \((1,2,1)\), and \((2,1,1)\), for example, span a
three-dimensional degenerate eigenspace.

Let \(H_z\in\mathbb R^{N_z\times N_z}\) be the third one-dimensional
finite-difference Hamiltonian. Flatten with \(x\) fastest, then \(y\), then
\(z\). The three-dimensional Kronecker sum is
\begin{align}
  H_{3,h}
  ={}&I_{N_z}\otimes I_{N_y}\otimes H_x\\
     &+I_{N_z}\otimes H_y\otimes I_{N_x}\\
     &+H_z\otimes I_{N_y}\otimes I_{N_x}.
\end{align}
It has shape
\begin{equation}
  H_{3,h}
  \in\mathbb R^{N_xN_yN_z\times N_xN_yN_z}.
\end{equation}
The associated product eigenvectors and eigenvalues are
\begin{equation}
  \mathbf v_{n_z}^{(z)}\otimes
  \mathbf v_{n_y}^{(y)}\otimes
  \mathbf v_{n_x}^{(x)},
  \qquad
  \varepsilon_{n_x}^{(x)}
  +\varepsilon_{n_y}^{(y)}
  +\varepsilon_{n_z}^{(z)}.
\end{equation}
Each identity factor means that the corresponding coordinate index is preserved
while another one-dimensional Hamiltonian acts.

\section{Sparse computational construction}
\index{sparse matrix!Kronecker construction}

A multidimensional finite-difference Hamiltonian is sparse. Each grid point
couples only to itself and to a bounded set of neighbors. The Kronecker
construction should preserve that sparsity rather than form a dense
\((N_xN_yN_z)\times(N_xN_yN_z)\) array.

\begin{pythonexample}{assembling sparse two- and three-dimensional boxes}
import numpy as np
from scipy.sparse import diags, eye, kron


def dirichlet_kinetic(point_count: int, length: float):
    # Use interior points; boundary values are fixed to zero and not stored.
    spacing = length / (point_count + 1)
    off_diagonal = -np.ones(point_count - 1, dtype=np.float64)
    diagonal = 2.0 * np.ones(point_count, dtype=np.float64)

    # Set hbar = mass = 1 for this dimensionless pedagogical example.
    core = diags(
        (off_diagonal, diagonal, off_diagonal),
        offsets=(-1, 0, 1),
        format="csr",
    )
    return core / (2.0 * spacing**2)


# Declare each one-dimensional factor and its identity explicitly.
Nx, Ny, Nz = 4, 3, 2
Hx = dirichlet_kinetic(Nx, length=1.0)
Hy = dirichlet_kinetic(Ny, length=1.5)
Hz = dirichlet_kinetic(Nz, length=2.0)
Ix = eye(Nx, format="csr")
Iy = eye(Ny, format="csr")
Iz = eye(Nz, format="csr")

# C-order flattening makes x the fastest coordinate.
H2 = kron(Iy, Hx, format="csr") + kron(Hy, Ix, format="csr")
H3 = (
    kron(kron(Iz, Iy, format="csr"), Hx, format="csr")
    + kron(kron(Iz, Hy, format="csr"), Ix, format="csr")
    + kron(kron(Hz, Iy, format="csr"), Ix, format="csr")
)

# Check shapes without converting either sparse operator to a dense array.
assert H2.shape == (Nx * Ny, Nx * Ny)
assert H3.shape == (Nx * Ny * Nz, Nx * Ny * Nz)

# Preserve sparsity in the Hermiticity diagnostic.
residual = H3 - H3.conjugate().transpose()
residual.eliminate_zeros()
assert residual.nnz == 0
print("2D shape and nonzeros:", H2.shape, H2.nnz)
print("3D shape and nonzeros:", H3.shape, H3.nnz)
\end{pythonexample}

The example constructs synthetic dimensionless operators. It checks assembly
shape and exact sparse Hermiticity only. It does not establish refinement,
continuum convergence, or a physical material result.

\documentcallout{Notebook to do}{Create
\path{examples/tutorials/research-monograph/foundations/particle_in_box_dimensions.ipynb}
with explicit imports, explanatory inline comments, hand-checked Kronecker
products, declared flattening order, product-state eigenpairs, square/cubic
degeneracies, and sparse-preserving two- and three-dimensional diagnostics.
Keep it unexecuted in the repository.}

\section{What changes outside a rectangular Cartesian box}

The simple Kronecker sum depends on separability. It works here because

\begin{itemize}
  \item the domain is a Cartesian product of intervals;
  \item the mass is a constant scalar;
  \item the boundary conditions separate by coordinate;
  \item the potential is zero, or more generally is a sum of one-coordinate
        terms; and
  \item the coordinate metric is diagonal and constant.
\end{itemize}

A coupled potential \(V(x,y)\), a nonrectangular boundary, a spatially varying
mass, or nonorthogonal coordinates can destroy the simple sum of independent
one-dimensional operators. Chapter~\ref{ch:training-periodic-geometry}
introduces the metric needed for nonorthogonal cells, and
Chapter~\ref{ch:training-bloch-finite-differences} shows why mixed derivatives
then add diagonal-neighbor stencil terms that are absent from the rectangular
box matrix above.

\section{Common mistakes}

\begin{enumerate}
  \item Writing \(H_y\otimes H_x\) when the separable Hamiltonian requires the
        Kronecker sum \(I_y\otimes H_x+H_y\otimes I_x\).
  \item Omitting identity factors and therefore changing the represented space.
  \item Reversing Kronecker factors without changing the flattening convention.
  \item Treating \(\mathbf v_y\otimes\mathbf v_x\) and
        \(\mathbf v_x\otimes\mathbf v_y\) as identical coordinate columns.
  \item Assuming every repeated numerical eigenvalue is an exact physical
        degeneracy without a tolerance and symmetry argument.
  \item Densifying a sparse multidimensional Hamiltonian merely to inspect or
        test it.
  \item Calling a finite-difference matrix the continuum Hamiltonian rather
        than a finite representation of a declared discretization.
\end{enumerate}

\section{Exercises}

\begin{enumerate}
  \item Compute \(A\otimes B\) by blocks for two chosen \(2\times2\) matrices,
        and verify its action on one vector product.
  \item Prove
        \((A\otimes B)(\mathbf x\otimes\mathbf y)
        =(A\mathbf x)\otimes(B\mathbf y)\)
        directly from components.
  \item Starting from one-dimensional box eigenfunctions, verify the
        normalization of \(\psi_{n_x,n_y}(x,y)\).
  \item Derive the additive two-dimensional energy from separation of
        variables.
  \item Reconstruct the displayed \(4\times4\) matrix
        \(I_2\otimes K+K\otimes I_2\) one term at a time.
  \item For a cubic box, list the degeneracy due to permutations of
        \((1,2,3)\), and distinguish those states from possible accidental
        degeneracies involving different triples.
  \item Show that the three-dimensional product vector is an eigenvector of
        \(H_{3,h}\) and that its eigenvalue is the sum of the three
        one-dimensional eigenvalues.
  \item Describe the permutation needed if a code changes from \(x\)-fastest
        to \(z\)-fastest flattening.
\end{enumerate}

\section{Evidence boundary}

The continuum product formulas are exact for the stated rectangular Dirichlet
model. The finite Kronecker sums are exact identities for the chosen
one-dimensional matrices and ordering. Neither statement establishes numerical
convergence at an untested grid, adequacy for a nonseparable model, or relevance
to a semiconductor material. Those stronger claims require separate numerical
and scientific evidence.
